\begin{table}[H]
\scriptsize
\setlength{\tabcolsep}{2.5pt}
\renewcommand{\arraystretch}{0.95}

\caption{\label{tab:country-heterogeneity-counterfactual-gap}Observed and no-policy counterfactual inflation inequality, 2020--2024}
\centering
\resizebox{\textwidth}{!}{%
\begin{tabular}{lrrrrrrr}
\toprule
Country & \shortstack{Q1 cumulative\\ inflation} & \shortstack{Q5 cumulative\\ inflation} & \shortstack{Inflation\\ gap} & \shortstack{Q1 counterfactual\\ cumulative inflation} & \shortstack{Q5 counterfactual\\ cumulative inflation} & \shortstack{Counterfactual\\ inflation gap} & \shortstack{Effect of fiscal\\ policies}\\
\midrule
Austria & 23.4 & 24.0 & -0.7 & 23.4 & 24.0 & -0.7 & 0.0\\
Cyprus & 17.8 & 17.5 & 0.2 & 17.8 & 17.5 & 0.3 & -0.0\\
Estonia & 46.6 & 39.8 & 6.9 & 46.6 & 39.8 & 6.9 & -0.0\\
Finland & 13.6 & 15.9 & -2.2 & 13.6 & 15.9 & -2.2 & 0.0\\
France & 17.8 & 17.9 & -0.1 & 17.8 & 17.9 & -0.1 & 0.0\\
\addlinespace
Germany & 20.0 & 22.5 & -2.5 & 20.1 & 22.5 & -2.5 & -0.0\\
Greece & 17.9 & 17.6 & 0.3 & 17.9 & 17.6 & 0.3 & 0.0\\
Ireland & 21.4 & 16.0 & 5.5 & 21.7 & 16.1 & 5.6 & -0.1\\
Italy & 19.8 & 19.8 & -0.0 & 20.5 & 20.7 & -0.1 & 0.1\\
Latvia & 38.1 & 32.1 & 6.1 & 38.1 & 32.1 & 6.1 & 0.0\\
\addlinespace
Lithuania & 38.2 & 35.5 & 2.7 & 38.2 & 35.5 & 2.7 & -0.0\\
Luxembourg & 16.6 & 17.8 & -1.1 & 16.6 & 17.8 & -1.1 & -0.0\\
Netherlands & 22.3 & 23.6 & -1.3 & 22.2 & 23.5 & -1.3 & 0.0\\
Portugal & 17.9 & 18.0 & -0.1 & 18.1 & 18.2 & -0.1 & -0.0\\
Slovakia & 31.2 & 32.3 & -1.1 & 31.2 & 32.3 & -1.1 & -0.0\\
\addlinespace
Spain & 18.3 & 18.5 & -0.2 & 18.8 & 18.8 & -0.0 & -0.1\\
\textbf{Euro area} & \textbf{19.8} & \textbf{20.6} & \textbf{-0.7} & \textbf{20.0} & \textbf{20.8} & \textbf{-0.7} & \textbf{0.0}\\
\bottomrule
\end{tabular}%
}
\begin{flushleft}
\footnotesize{\emph{Notes:} Cumulative inflation compares average income-quintile price indices in 2020 and 2024. Inflation gaps are Quintile 1 minus Quintile 5. The first three columns reproduce Table~\ref{tab:country-heterogeneity-cumulative-gap}. Counterfactual indices use the same no-policy COICOP component paths reported in Appendix~\ref{sec:app-counterfactuals}; observed 2019 HICP data are used only to chain the January 2020 starting values. They replace observed prices for the affected components with estimated no-policy paths and otherwise retain observed prices; countries without a simulated measure are treated as unchanged. Household-group and euro-area aggregation use the same RAS expenditure and HICP country weights as the observed indices. The last column is the observed gap minus the counterfactual gap. All numeric columns are expressed in percentage points. \emph{Sources:} Eurostat; national statistical institutes; authors' calculations.}
\end{flushleft}
\end{table}
